\documentclass[border=5mm]{standalone}
% ==============================================================================
% SETUP.TEX - CONFIGURACIÓN GLOBAL DEL PROYECTO
% ==============================================================================
% Este archivo centraliza todos los paquetes y estilos.
% Debe incluirse en el preámbulo del libro principal y de los archivos standalone.
% \PassOptionsToPackage{gray}{xcolor}
% --- 1. CODIFICACIÓN E IDIOMA ---
% fix-cm habilita las fuentes Computer Modern en tamaños arbitrarios (por debajo
% de 5 pt, entre otros). Debe cargarse antes que fontenc.
\usepackage{fix-cm}
\usepackage[utf8]{inputenc}
\usepackage[T1]{fontenc}
\usepackage[spanish, es-tabla]{babel}  
\usepackage{csquotes} % Recomendado para babel y citas

\usepackage{import}
% mode=image: Usa el PDF pre-compilado, nunca intenta recompilar.
% Debes compilar cada figura manualmente antes de compilar el libro.
\usepackage[mode=image, subpreambles=false]{standalone}

% --- 2. MATEMÁTICAS Y CIENCIAS ---
\usepackage{amsmath, amssymb, amsfonts, mathtools}
\usepackage{enumitem}

% LaTeX trae \DeclareMathSizes{5}{5}{5}{5}, así que a 5 pt (\tiny) los subíndices
% salen del mismo tamaño que la base: en las etiquetas de figuras el M de
% $\omega_M$ queda desproporcionado. Se restablece la proporción habitual.
\DeclareMathSizes{5}{5}{3.7}{3}

% --- 3. DISEÑO Y GEOMETRÍA --- 
% Unificamos los márgenes para todo el documento
\usepackage[left=2.5cm,right=2.5cm,top=3cm,bottom=3cm]{geometry}
\makeatletter
\ifstandalone % Evita que geometry dañe el recorte de las figuras bajándolas 3cm
  \@ifclasswith{standalone}{crop=false}{
    % Es un capítulo/sección: Dejamos que geometry actúe.
  }{
    % Es una figura: Neutralizamos geometry para que no las recorte mal.
    \geometry{pass}
  }
\fi
\makeatother
\usepackage{parskip} % Espaciado entre párrafos sin sangría
\usepackage{float}   % Para usar [H] en figura

% Ajustes de flotantes para evitar espacios verticales exagerados
\renewcommand{\topfraction}{0.95}      % Permite que las figuras ocupen hasta el 95% superior
\renewcommand{\bottomfraction}{0.95}   % Permite que las figuras ocupen hasta el 95% inferior
\renewcommand{\textfraction}{0.05}     % Permite hojas con tan solo un 5% de texto
\renewcommand{\floatpagefraction}{0.8} % Requiere que una página de solo figuras esté 80% llena
\setcounter{topnumber}{4}
\setcounter{bottomnumber}{4}
\setcounter{totalnumber}{10}

% --- 4. GRÁFICOS Y TABLAS ---
\usepackage{graphicx}
\usepackage{booktabs} % Tablas profesionales
\usepackage{caption}  % Control de captions y \captionof
\usepackage{tikz}
\usetikzlibrary{arrows.meta, calc, angles, quotes, babel, backgrounds}
\usepackage{pgfplots}
\usepgfplotslibrary{groupplots}
\usepgfplotslibrary{external}
\usepgfplotslibrary{patchplots}
\pgfplotsset{compat=1.18}
\usepackage[american, straightvoltages]{circuitikz}

% --- 5. COLORES Y CAJAS ---

\usepackage[table]{xcolor}
\usepackage{tcolorbox}

% Definición de colores corporativos/estilo
\definecolor{utnblue}{RGB}{0, 51, 102}
\definecolor{codegreen}{rgb}{0,0.6,0}
\definecolor{codegray}{rgb}{0.5,0.5,0.5}
\definecolor{codepurple}{rgb}{0.58,0,0.82}
\definecolor{backcolour}{rgb}{0.96,0.96,0.96}

% --- 6. ENTORNOS PERSONALIZADOS (CAJAS) ---

% Caja "Resultado" (Azul Corporativo) — Definiciones, fórmulas clave, resultados importantes.
% Uso: \begin{resultado}[Fórmula de Euler] ... \end{resultado}
\newtcolorbox{resultado}[1][]{
    colback=utnblue!5!white,
    colframe=utnblue,
    title=\textbf{#1},
    fonttitle=\bfseries,
    sharp corners,
    boxrule=0.5mm
}

% Caja "Nota" (Gris) — Advertencias, aplicaciones, contexto secundario.
% Uso: \begin{nota}[Cuidado con la calculadora] ... \end{nota}
% Uso: \begin{nota} ... \end{nota}  → título por defecto "Nota"
\newtcolorbox{nota}[1][Nota]{
    colback=codegray!5!white,
    colframe=codegray,
    title=\textbf{#1},
    fonttitle=\bfseries,
    sharp corners,
    boxrule=0.5mm
}

% Caja inline para resaltar un valor y enlazarlo con su otra aparición en una tabla
% o en una cuenta: mismo color, mismo valor.
% Uso: \marcavalor{$4$}  o  \marcavalor[codegreen]{$4$}
\newtcbox{\marcavalor}[1][utnblue]{
    on line,
    boxsep=1pt, left=2pt, right=2pt, top=1pt, bottom=1pt,
    colback=#1!10!white,
    colframe=#1,
    boxrule=0.4pt,
    arc=2pt
}

% --- 7. CÓDIGO FUENTE (LISTINGS) ---
\usepackage{listings}

\lstdefinestyle{miestilo}{
    backgroundcolor=\color{backcolour},   
    commentstyle=\color{codegreen},
    keywordstyle=\color{magenta},
    numberstyle=\tiny\color{codegray},
    stringstyle=\color{codepurple},
    basicstyle=\ttfamily\footnotesize,
    breakatwhitespace=false,         
    breaklines=true,                 
    captionpos=b,                    
    keepspaces=true,                 
    numbers=left,                    
    numbersep=5pt,                  
    showspaces=false,                
    showstringspaces=false,
    showtabs=false,                  
    tabsize=2,
    frame=single,                    
    rulecolor=\color{codegray}
}

\lstset{style=miestilo}
% Soporte para tildes y ñ en bloques de código
\lstset{literate={ñ}{{\~n}}1 {á}{{\'a}}1 {é}{{\'e}}1 {í}{{\'i}}1 {ó}{{\'o}}1 {ú}{{\'u}}1}

% Operadores matemáticos personalizados

% Vector en sentido discreto (lista finita de coordenadas): negrita + flecha.
% Uso: \vect{v}, \vect{e}_k, \vect{0}.
\newcommand{\vect}[1]{\vec{\mathbf{#1}}}

% Fecha de versión y comando para capítulos standalone
\providecommand{\verdate}{\the\year.\ifnum\month<10 0\fi\the\month.\ifnum\day<10 0\fi\the\day}
\newcommand{\chapterversion}{%
    \vspace{-1.5cm}\begin{flushright}\small\textit{\color{codegray}Versión~\verdate}\end{flushright}\vspace{0.5cm}%
}

% --- 8. HIPERVÍNCULOS (DEBE SER EL ÚLTIMO PAQUETE) ---
\usepackage[colorlinks=true, linkcolor=utnblue, urlcolor=utnblue, citecolor=utnblue]{hyperref}

% --- 9. REFERENCIAS CRUZADAS ENTRE CAPÍTULOS STANDALONE ---
% Cuando se compila un capítulo standalone, xr-hyper lee los .aux de los
% otros capítulos (si existen) para resolver \ref{} a labels externos.
% En la compilación del libro completo este bloque no se activa.
\ifstandalone
  \usepackage{xr-hyper}
  \makeatletter
  \newcommand{\xrchap}[1]{%
    \edef\xrc@self{\detokenize{Capitulo#1}}%
    \edef\xrc@job{\jobname}%
    \ifx\xrc@self\xrc@job\else
      \IfFileExists{../#1capitulo/Capitulo#1.aux}%
        {\externaldocument{../#1capitulo/Capitulo#1}}{}%
    \fi
  }
  \makeatother
  \xrchap{01}\xrchap{02}\xrchap{03}\xrchap{04}
  \xrchap{05}\xrchap{06}\xrchap{07}\xrchap{08}
  \xrchap{09}\xrchap{10}\xrchap{11}\xrchap{12}
  \xrchap{13}
\fi
% \PassOptionsToPackage{gray}{xcolor}

% Procedimiento de invertir y desplazar paso a paso para
%   x(τ) = e^{-aτ} u(τ)        (a = 1)
%   h(τ) = e^{-bτ} u(τ)        (b = 2)
% evaluado en t₀ = 1.
%
% Cuatro paneles (2x2), uno por etapa de la receta:
%   (a) Representar:           x(τ) y h(τ) sobre el eje τ
%   (b) Invertir h:            h(-τ) = e^{bτ} u(-τ)
%   (c) Desplazar:             h(t₀-τ), peak en τ=t₀
%   (d) Multiplicar e integrar: producto x(τ)·h(t₀-τ) y área = y(t₀)

\begin{document}
\begin{tikzpicture}[>=Latex, font=\small]

\colorlet{xcol}{utnblue}
\colorlet{hcol}{red!80!black}
\colorlet{prodcol}{black}

\pgfmathsetmacro{\tcero}{1.0}
\pgfmathsetmacro{\aval}{1.0}
\pgfmathsetmacro{\bval}{2.0}

% Mismo rango y tamaño en los cuatro paneles (igual escala de comparación)
\pgfmathsetmacro{\xmn}{-1.6}
\pgfmathsetmacro{\xmx}{3.0}
\pgfmathsetmacro{\ymn}{-0.18}
\pgfmathsetmacro{\ymx}{1.25}

% =====================================================================
% PANEL (a): Representar x(τ) y h(τ)
% =====================================================================
\begin{axis}[
    name=panA,
    width=8cm, height=5cm,
    axis lines=center,
    xlabel={$\tau$},
    ylabel={},
    xlabel style={anchor=north west, at={(axis description cs:1,0)}},
    xmin=\xmn, xmax=\xmx,
    ymin=\ymn, ymax=\ymx,
    clip=true,
    axis line style={->, thick},
    xtick={0},
    xticklabels={$0$},
    ytick={0,1},
    every tick label/.style={font=\tiny},
    title={\footnotesize (a)\; Representar},
]
    % x(τ) = e^{-aτ} u(τ)
    \addplot[xcol, very thick, domain=\xmn:0, samples=2] {0};
    \addplot[xcol, very thick, domain=0:\xmx, samples=120]
        {exp(-\aval*x)};
    \draw[xcol, very thick] (axis cs:0,0) -- (axis cs:0,1);
    \fill[xcol]  (axis cs:0,1) circle (2pt);
    \node[xcol, font=\tiny, above right] at (axis cs:1.3,0.27) {$x(\tau)$};

    % h(τ) = e^{-bτ} u(τ)
    \addplot[hcol, very thick, domain=\xmn:0, samples=2] {0};
    \addplot[hcol, very thick, domain=0:\xmx, samples=120]
        {exp(-\bval*x)};
    \draw[hcol, very thick] (axis cs:0,0) -- (axis cs:0,1);
    \fill[hcol]  (axis cs:0,1) circle (2pt);
    \node[hcol, font=\tiny, below right] at (axis cs:0.7,0.18) {$h(\tau)$};
\end{axis}

% =====================================================================
% PANEL (b): Invertir h
% =====================================================================
\begin{axis}[
    name=panB,
    at={(panA.east)}, anchor=west, xshift=0.8cm,
    width=8cm, height=5cm,
    axis lines=center,
    xlabel={$\tau$},
    ylabel={},
    xlabel style={anchor=north west, at={(axis description cs:1,0)}},
    xmin=\xmn, xmax=\xmx,
    ymin=\ymn, ymax=\ymx,
    clip=true,
    axis line style={->, thick},
    xtick={0},
    xticklabels={$0$},
    ytick={0,1},
    every tick label/.style={font=\tiny},
    title={\footnotesize (b)\; Invertir $h$},
]
    % x(τ) -- se mantiene como referencia
    \addplot[xcol, very thick, domain=\xmn:0, samples=2] {0};
    \addplot[xcol, very thick, domain=0:\xmx, samples=120]
        {exp(-\aval*x)};
    \draw[xcol, very thick] (axis cs:0,0) -- (axis cs:0,1);
    \fill[xcol]  (axis cs:0,1) circle (2pt);
    \node[xcol, font=\tiny, above right] at (axis cs:1.3,0.27) {$x(\tau)$};

    % h(-τ) = e^{bτ} u(-τ): cero a la derecha del origen, crece hasta 1 en τ=0
    \addplot[hcol, very thick, domain=\xmn:0, samples=120]
        {exp(\bval*x)};
    \draw[hcol, very thick] (axis cs:0,0) -- (axis cs:0,1);
    \fill[hcol]  (axis cs:0,1) circle (2pt);
    % a la derecha del origen el valor es 0
    \addplot[hcol, very thick, domain=0:\xmx, samples=2] {0};
    \node[hcol, font=\tiny] at (axis cs:-1.0,0.55) {$h(-\tau)$};

    % flecha que sugiere la reflexión
    \draw[->, gray!60, thin]
        (axis cs:0.4,0.55) to[out=160, in=20] (axis cs:-0.4,0.55);
\end{axis}

% =====================================================================
% PANEL (c): Desplazar
% =====================================================================
\begin{axis}[
    name=panC,
    at={(panA.south)}, anchor=north, yshift=-1.0cm,
    width=8cm, height=5cm,
    axis lines=center,
    xlabel={$\tau$},
    ylabel={},
    xlabel style={anchor=north west, at={(axis description cs:1,0)}},
    xmin=\xmn, xmax=\xmx,
    ymin=\ymn, ymax=\ymx,
    clip=true,
    axis line style={->, thick},
    xtick={0,\tcero},
    xticklabels={$0$, $t_0$},
    ytick={0,1},
    every tick label/.style={font=\tiny},
    title={\footnotesize (c)\; Desplazar},
]
    % x(τ)
    \addplot[xcol, very thick, domain=\xmn:0, samples=2] {0};
    \addplot[xcol, very thick, domain=0:\xmx, samples=120]
        {exp(-\aval*x)};
    \draw[xcol, very thick] (axis cs:0,0) -- (axis cs:0,1);
    \fill[xcol]  (axis cs:0,1) circle (2pt);
    \node[xcol, font=\tiny, above right] at (axis cs:1.5,0.22) {$x(\tau)$};

    % h(t₀-τ) = e^{-b(t₀-τ)} para τ ≤ t₀
    \addplot[hcol, very thick, domain=\xmn:\tcero, samples=120]
        {exp(\bval*(x-\tcero))};
    \draw[hcol, very thick] (axis cs:\tcero,0) -- (axis cs:\tcero,1);
    \fill[hcol]  (axis cs:\tcero,1) circle (2pt);
    \fill[white] (axis cs:\tcero,0) circle (2pt);
    \draw[hcol]  (axis cs:\tcero,0) circle (2pt);
    \addplot[hcol, very thick, domain=\tcero:\xmx, samples=2] {0};
    \node[hcol, font=\tiny] at (axis cs:-0.35,0.35) {$h(t_0\!-\!\tau)$};

    % marcador t₀
    \draw[dashed, gray!60, thin] (axis cs:\tcero,-0.16) -- (axis cs:\tcero,1.1);
\end{axis}

% =====================================================================
% PANEL (d): Multiplicar e integrar
% =====================================================================
\begin{axis}[
    name=panD,
    at={(panB.south)}, anchor=north, yshift=-1.0cm,
    width=8cm, height=5cm,
    axis lines=center,
    xlabel={$\tau$},
    ylabel={},
    xlabel style={anchor=north west, at={(axis description cs:1,0)}},
    xmin=\xmn, xmax=\xmx,
    ymin=\ymn, ymax=\ymx,
    clip=true,
    axis line style={->, thick},
    xtick={0,\tcero},
    xticklabels={$0$, $t_0$},
    ytick={0,1},
    every tick label/.style={font=\tiny},
    title={\footnotesize (d)\; Multiplicar e integrar},
]
    % Área sombreada = integral del producto sobre [0, t₀].
    % producto en [0, t₀]: x(τ)·h(t₀-τ) = e^{-aτ}·e^{-b(t₀-τ)}
    %                                  = e^{-b t₀}·e^{(b-a)τ}
    \addplot[draw=none, fill=gray!28, domain=0:\tcero, samples=120]
        {exp(-\bval*\tcero) * exp((\bval-\aval)*x)} \closedcycle;

    % x(τ) en versión tenue (referencia)
    \addplot[xcol!35, very thick, domain=\xmn:0, samples=2] {0};
    \addplot[xcol!35, very thick, domain=0:\xmx, samples=120]
        {exp(-\aval*x)};
    \draw[xcol!35, very thick] (axis cs:0,0) -- (axis cs:0,1);

    % h(t₀-τ) en versión tenue
    \addplot[hcol!35, very thick, domain=\xmn:\tcero, samples=120]
        {exp(\bval*(x-\tcero))};
    \draw[hcol!35, very thick] (axis cs:\tcero,0) -- (axis cs:\tcero,1);

    % Producto x(τ)·h(t₀-τ) en negro sobre [0, t₀]
    \addplot[prodcol, very thick, domain=0:\tcero, samples=120]
        {exp(-\bval*\tcero) * exp((\bval-\aval)*x)};
    \draw[prodcol, very thick]
        (axis cs:0,0) -- (axis cs:0,{exp(-\bval*\tcero)});
    \draw[prodcol, very thick]
        (axis cs:\tcero,0) -- (axis cs:\tcero,{exp(-\aval*\tcero)});

    % etiqueta del producto y del área
    \node[prodcol, font=\tiny, above]
        at (axis cs:\tcero,{exp(-\aval*\tcero)})
        {$x(\tau)\,h(t_0\!-\!\tau)$};
    \node[gray!60!black, font=\tiny]
        at (axis cs:0.5,0.10) {$y(t_0)$};

    % marcador t₀
    \draw[dashed, gray!60, thin] (axis cs:\tcero,-0.16) -- (axis cs:\tcero,1.1);
\end{axis}

\end{tikzpicture}
\end{document}
